\subsection{模上的若干运算}

设 A 与 B 是环，P 是一个 (A, B)-双模（II, §1, No.~14, 第 225 页）。我们将定义两个手续：一个从左 B-模过渡到左 A-模，另一个从左 A-模过渡到左 B-模。

3.1. 运算 \(\mathcal{T}\) . —— 设 V 是一个左 B-模。把左 A-模 \(\mathrm{P} \otimes_{\mathrm{B}} \mathrm{V}\)（II, §3, No.~4, 第 247 页）记作 \(\mathcal{T}(\mathrm{V})\)。\(\mathcal{T}(\mathrm{V})\) 上的作用律由公式

\[
a (p \otimes v) = (a p) \otimes v\tag{1}
\]

给出，其中 \(a \in \mathrm{A}\)、\(p \in \mathrm{P}\)、\(v \in \mathrm{V}\) 。

设 \(V'\) 是一个左 B-模。对每个从 V 到 \(V'\) 的 B-线性映射 g，映射 \(1_{\mathrm{P}} \otimes g\)（从 \(\mathcal{T}(\mathrm{V})\) 到 \(\mathcal{T}(\mathrm{V}')\)）是 A-线性的；我们把它记作 \(\mathcal{T}(g)\)。映射 \(g \mapsto \mathcal{T}(g)\)（从 \(\operatorname{Hom}_{\mathrm{B}}(\mathrm{V}, \mathrm{V}')\) 到 \(\operatorname{Hom}_{\mathrm{A}}(\mathcal{T}(\mathrm{V}), \mathcal{T}(\mathrm{V}'))\)）是 Z-线性的，且我们有

\[
\mathcal {T} (1 _ {\mathrm{V}}) = 1 _ {\mathcal {T} (\mathrm{V})}, \quad \mathcal {T} (g ^ {\prime} \circ g) = \mathcal {T} (g ^ {\prime}) \circ \mathcal {T} (g)\tag{2}
\]

（其中 V、\(V'\)、\(V''\) 是左 B-模，且 \(g : V \to V'\)、\(g' : V' \to V''\) 是 B-线性映射）。由于张量积与直和交换，若 V 是一族子模 \((\mathrm{V}_{i})_{i \in \mathrm{I}}\) 的直和，则我们可以把 A-模 \(\mathcal{T}(\mathrm{V})\) 与 \(\bigoplus_{i} \mathcal{T}(\mathrm{V}_{i})\) 等同起来。

3.2. 运算 \(\mathcal{H}\) . —— 设 M 是一个左 A-模。把左 B-模 \(\mathrm{Hom}_{\mathrm{A}}(\mathrm{P},\mathrm{M})\)（II, §1, No.~14, 第 225 页）记作 \(\mathcal{H}(\mathrm{M})\)。\(\mathcal{H}(\mathrm{M})\) 上的作用律由公式

\[
(b f) (p) = f (p b)\tag{3}
\]

给出，其中 \(b \in \mathrm{B}\)、\(f \in \operatorname{Hom}_{\mathrm{A}}(\mathrm{P}, \mathrm{M})\)、\(p \in \mathrm{P}\) 。

设 \(M'\) 是一个左 A-模。对每个从 M 到 \(M'\) 的 A-线性映射 g，映射 \(\mathrm{Hom}_{\mathrm{A}}(1_{\mathrm{P}}, g)\)（从 \(\mathcal{H}(\mathrm{M})\) 到 \(\mathcal{H}(\mathrm{M}')\)）是 B-线性的；我们把它记作 \(\mathcal{H}(g)\)。映射 \(g \mapsto \mathcal{H}(g)\)（从 \(\mathrm{Hom}_{\mathrm{A}}(\mathrm{M}, \mathrm{M}')\) 到 \(\mathrm{Hom}_{\mathrm{B}}(\mathcal{H}(\mathrm{M}), \mathcal{H}(\mathrm{M}'))\)）是 Z-线性的，且我们有

\[
\mathcal {H} (1 _ {\mathrm{M}}) = 1 _ {\mathcal {H} (\mathrm{M})}, \quad \mathcal {H} (g ^ {\prime} \circ g) = \mathcal {H} (g ^ {\prime}) \circ \mathcal {H} (g)\tag{4}
\]

（其中 M、\(M'\)、\(M''\) 是左 A-模，且 \(g : M \to M'\)、\(g' : M' \to M''\) 是 A-线性映射）。再设 P 是一个有限生成 A-模；若 M 是一族子模 \((\mathrm{M}_{i})_{i}\) 的直和，则由 VIII，第 57 页，我们可以把 \(\mathcal{H}(\mathrm{M})\) 与 \(\bigoplus_{i} \mathcal{H}(\mathrm{M}_{i})\) 等同起来。

3.3. \(\mathcal{T}\) 与 \(\mathcal{H}\) 之间的关系 . —— 由 II, §4, No.~1, 第 268 页的命题 1，对每个左 A-模 M 与每个左 B-模 V，存在唯一的群同构

\[
\gamma : \mathrm{Hom} _ {\mathrm{A}} (\mathcal {T} (\mathrm{V}), \mathrm{M}) \longrightarrow \mathrm{Hom} _ {\mathrm{B}} (\mathrm{V}, \mathcal {H} (\mathrm{M}))\tag{5}
\]

它由关系式

\[
(\gamma (h) (v)) (p) = h (p \otimes v)\tag{6}
\]

刻画，其中 \(h \in \operatorname{Hom}_{\mathrm{A}}(\mathcal{T}(\mathrm{V}), \mathrm{M})\)、\(v \in V\)、\(p \in P\)。同构 \(\gamma\) 称为伴随同构。

设 M 是一个左 A-模。A-模 \(\mathcal{T}(\mathcal{H}(M))\) 就是 A-模 \(P \otimes_{B} \operatorname{Hom}_{A}(P, M)\)。把上述结果应用于 B-模 \(\mathcal{H}(M)\)，我们看出映射

\[
\alpha_ {\mathrm{M}} = \gamma^ {- 1} (\operatorname{Id} _ {\mathcal {H} (\mathrm{M})}): \mathcal {T} (\mathcal {H} (\mathrm{M})) \longrightarrow \mathrm{M}
\]

是满足

\[
\alpha_ {\mathrm{M}} (p \otimes f) = f (p)\tag{7}
\]

的唯一映射，其中 \(p \in \mathrm{P}\)、\(f \in \mathrm{Hom}_{\mathrm{A}}(\mathrm{P}, \mathrm{M})\)。我们称 \(\alpha_{\mathrm{M}}\) 为从 \(\mathcal{T}(\mathcal{H}(\mathrm{M}))\) 到 M 的典范 A-线性映射。对每个 A-线性映射 \(g: \mathrm{M} \to \mathrm{M}'\)，我们有一个交换图

\[
\begin{array}{c} \mathcal {T} (\mathcal {H} (M)) \xrightarrow {\alpha_ {M}} M \\ \Biggl \downarrow \mathcal {T} (\mathcal {H} (g)) \Biggl \downarrow g \\ \mathcal {T} (\mathcal {H} (M ^ {\prime})) \xrightarrow {\alpha_ {M ^ {\prime}}} M ^ {\prime}. \end{array}\tag{I}
\]

伴随同构的逆

\[
\gamma^ {- 1}: \mathrm{Hom} _ {\mathrm{B}} (\mathrm{V}, \mathcal {H} (\mathrm{M})) \longrightarrow \mathrm{Hom} _ {\mathrm{A}} (\mathcal {T} (\mathrm{V}), \mathrm{M})
\]

与映射 \(h \mapsto \alpha_{\mathrm{M}} \circ \mathcal{T}(h)\) 一致。事实上，由 (6) 与 (7)，我们有关系式

\[
\gamma^ {- 1} (h) (p \otimes v) = (h (v)) (p) = \alpha_ {\mathrm{M}} (p \otimes h (v)) = \alpha_ {\mathrm{M}} \circ \mathcal {T} (h) (p \otimes v)
\]

对一切 \(h\in \mathrm{Hom}_{\mathrm{B}}(\mathrm{V},\mathcal{H}(\mathrm{M}))\)、\(v\in \mathrm{V}\)、\(p\in \mathrm{P}\) 成立。

设 V 是一个 B-模。B-模 \(\mathcal{H}(\mathcal{T}(V))\) 就是 B-模 \(\mathrm{Hom}_{\mathrm{A}}(\mathrm{P},\mathrm{P}\otimes_{\mathrm{B}}\mathrm{V})\)。把 (5) 应用于 A-模 \(\mathcal{T}(V)\)，我们看出 B-线性映射 \(\beta_{\mathrm{V}}=\gamma(\mathrm{Id}_{\mathcal{T}(V)})\)（从 V 到 \(\mathcal{H}(\mathcal{T}(V))\)）由关系式

\[
\beta_ {\mathrm{V}} (v) (p) = p \otimes v\tag{8}
\]

刻画，其中 \(p \in P\)、\(v \in V\)。我们称 \(\beta_{V}\) 为从 V 到 \(\mathcal{H}(\mathcal{T}(V))\) 的典范 B-线性映射。对每个 B-模 \(V'\) 与每个 B-线性映射 \(g : V \to V'\)，我们有一个交换图

\[
\begin{array}{l} \mathrm{V} \xrightarrow {\beta_ {\mathrm{V}}} \mathcal {H} (\mathcal {T} (\mathrm{V})) \\ \Big \downarrow_ {g} \qquad \qquad \qquad \Big \downarrow_ {\mathcal {H} (\mathcal {T} (g))} \\ \mathrm{V} ^ {\prime} \xrightarrow {\beta_ {\mathrm{V} ^ {\prime}}} \mathcal {H} (\mathcal {T} (\mathrm{V} ^ {\prime})). \end{array}\tag{II}
\]

注意，伴随态射 (5) 与把 u 映为 \(\mathcal{H}(u)\circ\beta_{\mathrm{V}}\) 的映射一致。事实上，由关系式 (6) 与 (8)，我们导出等式

\[
(\gamma (u) (v)) (p) = u (p \otimes v) = u \circ (\beta_ {\mathrm{V}} (v)) (p)
\]

对一切 \(u \in \mathrm{Hom}_{\mathrm{A}}(\mathcal{T}(\mathrm{V}), \mathrm{M})\)、\(v \in \mathrm{V}\)、\(p \in \mathrm{P}\) 成立。

注. —— 1) 设 V 与 V' 是 B-模。伴随同构

\[
\gamma : \mathrm{Hom} _ {\mathrm{A}} (\mathcal {T} (\mathrm{V}), \mathcal {T} (\mathrm{V} ^ {\prime})) \longrightarrow \mathrm{Hom} _ {\mathrm{B}} (\mathrm{V}, \mathcal {H} (\mathcal {T} (\mathrm{V} ^ {\prime})))
\]

满足关系式 \(\gamma (\mathcal{T}(f)) = \beta_{\mathrm{V}^{\prime}}\circ f\)（对每个 \(f\in \mathrm{Hom}_{\mathrm{B}}(\mathrm{V},\mathrm{V}^{\prime})\)），因为

\[
(\gamma (\mathcal {T} (f)) (v)) (p) = \mathcal {T} (f) (p \otimes v) = p \otimes f (v) = (\beta_ {\mathrm{V} ^ {\prime}} \circ f) (v) (p).
\]

设 M 与 M' 是 A-模；伴随同构的逆

\[
\gamma^ {- 1}: \mathrm{Hom} _ {\mathrm{B}} (\mathcal {H} (\mathrm{M}), \mathcal {H} (\mathrm{M} ^ {\prime})) \longrightarrow \mathrm{Hom} _ {\mathrm{A}} (\mathcal {T} (\mathcal {H} (\mathrm{M})), \mathrm{M} ^ {\prime})
\]

满足关系式 \(\gamma^{-1}(\mathcal{H}(u)) = u \circ \alpha_{\mathrm{M}}\)（对每个 \(u \in \operatorname{Hom}_{\mathrm{B}}(\mathrm{M}, \mathrm{M}')\)）。事实上，我们有关系式

\[
\gamma^ {- 1} (\mathcal {H} (u)) (p \otimes v) = (\mathcal {H} (u) (v)) (p) = u (v (p)) = u \circ \alpha_ {\mathrm{M}} (p \otimes v)
\]

对一切 \(u \in \mathrm{Hom}_{\mathrm{B}}(\mathrm{M}, \mathrm{M}')\)、\(v \in \mathcal{H}(\mathrm{M})\)、\(p \in \mathrm{P}\) 成立。

\begin{enumerate}
\def\labelenumi{\arabic{enumi})}
\setcounter{enumi}{1}
\tightlist
\item
  设 \(\mathrm{M}\) 是一个左 A-模。B-线性映射
\end{enumerate}

\[
\beta_ {\mathcal {H} (\mathrm{M})}: \mathcal {H} (\mathrm{M}) \to \mathcal {H} (\mathcal {T} (\mathcal {H} (\mathrm{M})))
\]

\[
\mathcal {H} (\alpha_ {\mathrm{M}}): \mathcal {H} (\mathcal {T} (\mathcal {H} (\mathrm{M}))) \to \mathcal {H} (\mathrm{M})
\]

满足关系式 \(\mathcal{H}(\alpha_{\mathrm{M}})\circ\beta_{\mathcal{H}(\mathrm{M})}=1_{\mathcal{H}(\mathrm{M})}\)。一般而言它们不是双射。

设 V 是一个左 B-模。A-线性映射

\[
\mathcal {T} (\beta_ {\mathrm{V}}): \mathcal {T} (\mathrm{V}) \to \mathcal {T} (\mathcal {H} (\mathcal {T} (\mathrm{V}))) \quad \text { and } \quad \alpha_ {\mathcal {T} (\mathrm{V})}: \mathcal {T} (\mathcal {H} (\mathcal {T} (\mathrm{V}))) \to \mathcal {T} (\mathrm{V})
\]

满足关系式 \(\alpha_{\mathcal{T}(\mathrm{V})}\circ \mathcal{T}(\beta_{\mathrm{V}}) = 1_{\mathcal{T}(\mathrm{V})}\)。一般而言它们不是双射。

\begin{enumerate}
\def\labelenumi{\arabic{enumi})}
\setcounter{enumi}{2}
\tightlist
\item
  设 P 作为 A-模是有限生成的。设 M 是一族 A-模 \((\mathrm{M}_{i})_{i\in\mathrm{I}}\) 的直和。A-模 \(\mathcal{T}(\mathcal{H}(\mathrm{M}))\) 与 \(\bigoplus_{i}\mathcal{T}(\mathcal{H}(\mathrm{M}_{i}))\) 典范同构。当我们把它们等同起来时，\(\alpha_{\mathrm{M}}\) 就与 \(\bigoplus_{i}\alpha_{\mathrm{M}_{i}}\) 等同。同样，设 V 是一族 B-模 \((\mathrm{V}_j)_{j\in \mathrm{J}}\) 的直和。B-模 \(\mathcal{H}(\mathcal{T}(\mathrm{V}))\) 与 \(\bigoplus_{j}\mathcal{H}(\mathcal{T}(\mathrm{V}_{j}))\) 等同，且线性映射 \(\beta_{\mathrm{V}}\) 与 \(\bigoplus_{j}\beta_{\mathrm{V}_{j}}\) 等同。
\end{enumerate}

